\section{Framework and the quantitative step}\label{sec:framework}

\subsection{The weighted base operator}

Let $B=S^4\subset\R^5$, with round metric $g_0$ and round Riemannian
density $\mu$. Write a point of $B$ as $(x_1,x_2,y)$, where $y\in\R^3$,
and set $t=|y|^2$. A \emph{positive pair} consists of a smooth positive
definite symmetric contravariant tensor $A$ and a smooth function
$\rho>0$. Define a metric $G$ and a differential operator $L$ by
\begin{equation}\label{eq:weighted-operator}
 G^{-1}=\rho^{-1}A,\qquad
 Lf=\rho^{-1}\diver_\mu(A\,\dd f).
\end{equation}
Thus $-L$ is symmetric and nonnegative in $L^2(B,\rho\,\dd\mu)$, with
quadratic form
\[
 \mathfrak a(f,f)=\int_B A(\dd f,\dd f)\,\dd\mu .
\]
The tensor and density will be changed independently. Their conversion
to an ordinary metric occurs only at the end of the proof.

Fix a finite smooth atlas with compact coordinate subpatches covering
$B$; all $C^r$ norms of functions and tensors may be taken in that atlas.
For a pair, the norm means the sum of the two component norms.
Equivalent fixed choices change constants but do not change the
assertions. All pairs used below lie in one sufficiently small $C^1$
neighborhood of $(g_0^{-1},1)$, within the space of smooth positive
pairs. Accordingly $G$ and $G^{-1}$ have common $C^1$ bounds and common
ellipticity constants. Higher derivative bounds need only be finite for
each fixed input pair.

Gradients, Hessians, lengths, and orthogonality in envelope and phase
estimates refer to $G$ unless another metric is displayed. Integration
in the elliptic operator refers to $\mu$; integration in a coordinate
corrugation chart refers to Lebesgue measure $\dd x$.
For a real smooth function $h$, set
\begin{equation}\label{eq:hessian-test}
 \begin{split}
 p&=\nabla_G h,\qquad Q=(1+|p|_G^2)^{1/2},\\
 J_h(e)&=(\nabla_G^2h)(p/Q,p/Q)+(\nabla_G^2h)(e,e),
 \qquad |e|_G=1,\quad \dd h(e)=0 .
 \end{split}
\end{equation}
At $p=0$ every unit vector is an admissible test.
The positivity conditions imposed below on $J_h$ supply a gap between
the envelope Hessian and the real part of a complex phase Hessian.
Lemma~\ref{lem:phase-hessian} and Proposition~\ref{prop:packets} use
this gap in the construction and localization of the packets.
Complex inner products
and tensor evaluations in the phase equations are extended
bilinearly; estimates of the size of a complex vector or matrix use
the ordinary Hermitian norm.

We fix once and for all
\[
 \delta=10^{-6}.
\]
The numerical exponents below leave generous margins; we do not optimize
the final positive gain.
Positive constants written $c,C,C_r$ may change between occurrences.
Any additional dependencies are specified where their uniformity
matters. In particular, a finite approximation order is fixed before a
frequency tends to infinity.

\subsection{A round region and sign certificates}

For an integer $N>1$, let
\begin{equation}\label{eq:round-beam}
 k=\sqrt{N(N+3)},\qquad
 u_*=\operatorname{Re}(x_1+i x_2)^N .
\end{equation}
The Euclidean polynomial in \eqref{eq:round-beam} is homogeneous of
degree $N$ and harmonic in $\R^5$. The polar-coordinate formula
\[
 \Delta_{\R^5}
 =\partial_r^2+\frac4r\partial_r+\frac1{r^2}\Delta_{g_0}
\]
therefore gives $-\Delta_{g_0}u_*=k^2u_*$ on $B$.
Where $t<1$, the angular coordinate $\vartheta$ in the
$(x_1,x_2)$-plane also gives
\[
 u_*=(1-t)^{N/2}\cos(N\vartheta).
\]
The round beam will be retained in part of a region where the input
pair is round. This region serves two purposes: it supplies an exact
solution to which the packets can be glued, and it leaves a patch where
a later tensor perturbation can make the eigenvalue simple while
preserving that solution. The construction uses both properties before
passing to the limiting operator.

The envelope construction fixes a coordinate chart
$\Omega\subset\{t>1/2\}$ and an open region $V$ whose closure is compact
in $\Omega$. The initial corrugation cube is compactly contained in
$V$. Euclidean balls centered in $V$ with radii tending uniformly to
zero are consequently contained in $\Omega$.

\begin{definition}\label{def:certificate}
A \emph{sign certificate} for a continuous real function $u$ in
$\Omega$ is a finite family of pairs $(D_i^-,D_i^+)$ of congruent
closed parallel three-dimensional Euclidean disks such that:
\begin{enumerate}[label=\textup{(\roman*)}]
 \item each pair is the two ends of a straight closed cylinder
 compactly contained in a Euclidean coordinate ball;
 \item these coordinate balls are pairwise disjoint;
 \item $u$ has strictly opposite signs on $D_i^-$ and $D_i^+$.
\end{enumerate}
The \emph{size} of the certificate is
$\sum_i\Hh_{\mathrm{Euc}}^3(D_i^+)$.
\end{definition}

The signs in Definition~\ref{def:certificate} can be reversed on each
pair. Since the disks are compact and finite in number, a certificate
persists under sufficiently small uniform perturbations of $u$.
It is also preserved by multiplying $u$ by any nonzero real constant.
Every segment joining corresponding points of $D_i^-$ and $D_i^+$
contains a zero. This elementary projection property will provide the
Hausdorff-measure bound for the final metric.

\subsection{The one-step construction}

\begin{proposition}\label{prop:step}
There exist a $C^1$ neighborhood $\mathcal U$ of the round pair
$(g_0^{-1},1)$ among smooth positive pairs, a chart $\Omega$ as above,
and a number $\gamma>0$ with the following property.
Fix a smooth pair $(A,\rho)\in\mathcal U$, numbers
\[
 \frac14\le l_-<l_+\le\frac12,
\]
and suppose that $(A,\rho)=(g_0^{-1},1)$ on $\{t<l_+\}$.
For every integer $r\ge0$ and every $\eta>0$, there exist $c>0$ and
$N_0$ such that for every integer $N\ge N_0$, with $k$ as in
\eqref{eq:round-beam}, there are a smooth positive pair
$(\widehat A,\widehat\rho)\in\mathcal U$ and a nonzero real smooth
function $u$ satisfying:
\begin{enumerate}[label=\textup{(\roman*)}]
 \item $\|(\widehat A-A,\widehat\rho-\rho)\|_{C^r}<\eta$, and
 $(\widehat A,\widehat\rho)=(g_0^{-1},1)$ on $\{t<l_-\}$;
 \item for $\widehat L$ defined by \eqref{eq:weighted-operator},
 $-\widehat Lu=k^2u$, and $k^2$ is a simple eigenvalue of
 $-\widehat L$ on the base $B$;
 \item $u$ has a sign certificate in $\Omega$ of size at least
 $c k^{1+\gamma}$.
\end{enumerate}
The neighborhood $\mathcal U$, chart $\Omega$, and exponent $\gamma$
are common to all stages. The constant $c$ and threshold $N_0$ may
depend on the fixed input pair, $l_-,l_+,r,\eta$.
\end{proposition}

The proof of Proposition~\ref{prop:step} occupies the next three
sections. Membership of the output in $\mathcal U$ is ensured by
imposing an additional sufficiently small tolerance around the fixed
input pair. The conclusion permits arbitrarily small finite-norm
tolerances; it makes no claim of a uniform frequency threshold as
those tolerances or the collar width vary. The common exponent
$\gamma$ is the feature that allows a smooth limiting metric to retain
a fixed power of excess nodal growth.
